\section{Study Intervention}

The study intervention has two inseparable components delivered under a single
integrated protocol: an investigational drug, perioperative daraxonrasib
(RMC-6236), evaluated by a standard 3+3 dose-finding design under the
investigational new drug pathway; and an investigational device, the on-premises
large language model (LLM) directed eight-arm robotic pancreaticoduodenectomy
(Whipple) system, evaluated as an early-feasibility significant-risk Class II
collaborative platform under the investigational device exemption pathway. This
section describes both components, their administration, their preparation and
accountability, the measures taken to minimize bias, the compliance machinery for
each, and the permitted and prohibited concomitant therapy. The two components
are coupled at one decision point only: the perioperative pause-and-restart
advisory for the drug is informed by the operative and pathology results of the
device procedure ($\S$6.1.2).

\subsection{Study Intervention Administration}

\subsubsection{Study Intervention Description}

\paragraph{Drug component: daraxonrasib (RMC-6236).}
Daraxonrasib is an orally bioavailable, RAS(ON) multi-selective (pan-RAS)
small-molecule inhibitor that binds the active, guanosine triphosphate bound state
of mutant RAS in complex with cyclophilin A, suppressing downstream effector
signaling across the common KRAS oncoprotein variants found in pancreatic ductal
adenocarcinoma. The agent received FDA Breakthrough Therapy designation in
June 2025 for previously treated metastatic pancreatic ductal adenocarcinoma
\cite{revbreakthrough}, and its dose and exposure assumptions in this protocol are
anchored to the pan-KRAS RASolute 302 program, in which 300 mg once daily is the
established recommended Phase 2 dose \cite{rasolute302}. In this study daraxonrasib
is administered orally, once daily, in a perioperative schedule: a defined
pre-operative course, a protocol-mandated pause spanning the operative window, and
a postoperative restart whose timing is governed by the advisory described in
$\S$6.1.2. The drug is investigational in this perioperative, curative-intent
surgical context and for the combined drug-device use described here.

\paragraph{Device component: on-premises LLM-directed eight-arm Whipple system.}
The investigational device is an eight-arm parallel coelomic surgical platform in
which an on-premises repository-pinned LLM acts as a second-opinion advisory oracle
held strictly outside the robot vendor kinematic stack, never as an autonomous
controller. The platform carries 56 mechanical degrees of freedom ($8\times7$
per-arm degrees of freedom) and a 640-channel mixed-rate sensor stack (80 channels
per arm). Force channels sample at 100 kHz; all remaining channels sample at the
10 kHz command rate. Positioning accuracy is 0.05 mm root mean square at a peak tip
velocity of $1{,}200$ mm/s. The hard cyber-physical safety envelope comprises a
per-arm tip-force cap of $\leq 3$ N, a cumulative cross-arm tip-force cap of
$\leq 18$ N enforced at the heartbeat boundary, a 10 kHz heartbeat coordination bus
broadcasting a 64-byte frame on a $100\,\mu$s per-arm watchdog deadline, an
emergency arm park within $50\,\mu$s of a missed or out-of-parameter frame, and a
cross-arm emergency stop (E-stop) that halts the full platform in $\leq 3$ ms.
A software-independent hardware E-stop backs the software chain. In this Phase 1
protocol the device operates only at the clinician-controlled and supervised
semiautonomous autonomy levels under continuous human oversight with immediate
manual override; full autonomy is prohibited consistent with 21 CFR
$\S$312.21(e). The advisory control loop and the platform architecture are shown
in Figure 6 and Figure 7.

\begin{mermaidfig}
\node[mmin]  (sens) at (0,0)      {640 sensor\\channels\\100 kHz force /\\10 kHz other};
\node[mmstep](map)  at (4.2,0)    {Sensor to x,y,z\\Cartesian command\\frame};
\node[mmgoal](llm)  at (8.4,0)    {On-premises LLM\\advisory, hash-pinned\\to commit + seed};
\node[mmstep](gate) at (8.4,-2.6) {Safety gate\\vessel zones + force\\caps; 10 kHz\\heartbeat,\\$100\,\mu$s watchdog};
\node[mmin]  (vend) at (4.2,-2.6) {Robot vendor\\kinematic stack\\(isolated from LLM)};
\node[mmgoal](act)  at (0,-2.6)   {8-arm actuators\\$\leq 3$ N/arm,\\$\leq 18$ N cumulative};
\node[mmdark](aud)  at (0,-5.0)   {Hash-chained\\audit trail\\21 CFR part 11};
\draw[mmedge] (sens) -- (map);
\draw[mmedge] (map) -- (llm);
\draw[mmedge] (llm) -- (gate);
\draw[mmedge] (gate) -- (vend);
\draw[mmedge] (vend) -- (act);
\draw[mmedge] (act.south) -- (aud.north);
\draw[mmedge] (act.west) -- ++(-0.5,0) |- ($(sens.west)+(-0.5,0)$) -- (sens.west);
\draw[mmedged] (gate.south) -- ++(0,-0.5) -| (act.north);
\end{mermaidfig}
\figcaption{Figure 6. On-premises LLM advisory control loop. Sensor data is mapped
to a Cartesian command frame; the LLM issues hash-pinned advisory commands from
outside the vendor kinematic stack (second-opinion oracle isolation); and a safety
gate enforces the vessel zones, the force caps, the 10 kHz heartbeat, the
$100\,\mu$s watchdog, and the $\leq 3$ ms E-stop before any actuator moves. The
actuator telemetry returns to the sensor stack and every step is written to a
21 CFR part 11 hash-chained audit trail.}

\begin{mermaidfig}
\node[mmgoal](hub) at (4.2,0)     {10 kHz heartbeat\\coordination bus\\64-byte frame,\\$100\,\mu$s deadline};
\node[mmstep](a1)  at (0,-2.4)    {Arm 1 / Arm 2\\hybrid u-w-p\\dissection (P1-P8)};
\node[mmstep](a3)  at (4.2,-2.4)  {Arm 3 bipolar +\\retractor; Arm 4\\NIR + ICG (P1-P8)};
\node[mmstep](a5)  at (8.4,-2.4)  {Arm 5 anastomosis\\(P5-P7); Arm 6 coag\\(P2, P5-P8)};
\node[mmstep](a7)  at (8.4,-4.8)  {Arm 7 suction +\\irrigation (P1-P8)};
\node[mmstep](a8)  at (4.2,-4.8)  {Arm 8 imaging +\\drain (P1, P4, P8)};
\node[mmin]  (spec)at (0,-4.8)    {56 DOF ($8\times7$)\\640 channels (80/arm)\\0.05 mm RMS at\\$1{,}200$ mm/s};
\draw[mmedge] (hub) -- (a1);
\draw[mmedge] (hub) -- (a3);
\draw[mmedge] (hub) -- (a5);
\draw[mmedge] (a5) -- (a7);
\draw[mmedge] (a3) -- (a8);
\draw[mmedged] (a1) -- (spec);
\end{mermaidfig}
\figcaption{Figure 7. Eight-arm platform architecture. Eight cooperating arms
(56 degrees of freedom, 640 sensor channels at 80 per arm, 0.05 mm root mean square
positioning at $1{,}200$ mm/s) are coordinated by a 10 kHz heartbeat bus, with each
arm's tool assignment and operative-phase coverage driving the Schedule of
Activities.}

The per-arm tool assignment crossed with the eight operative phases (P1 through P8)
is given in Table 6, and ten representative channels drawn from the 640-channel
inventory are given in Table 7.

\begin{table}[H]
\caption{Per-arm tool assignment and operative-phase coverage for the eight-arm
LLM-directed Whipple platform.}
\label{tab:sec06-arms}
\centering
\begin{tabularx}{\textwidth}{L{1.2cm} L{4.4cm} L{4.0cm} Y}
\toprule
\textbf{Arm} & \textbf{Primary tool} & \textbf{Phase coverage} & \textbf{Sensor channels (10 of 80)} \\
\midrule
Arm 1 & Hybrid u-w-p end-effector & P1 to P8 (active dissect) & Tip force, joint torque, ID \\
Arm 2 & Hybrid u-w-p end-effector & P1 to P8 (active dissect) & Tip force, joint torque, ID \\
Arm 3 & Bipolar + retractor & P1 to P8 (vessel control) & Tip force, current, force \\
Arm 4 & NIR + ICG imaging probe & P1 to P8 (imaging) & ICG fluorescence, NIR depth \\
Arm 5 & Anastomosis probe & P5, P6, P7 (idle others) & Ring tension, ring strain \\
Arm 6 & Bipolar coag + cautery & P2, P5, P6, P7, P8 & Tip force, RF power \\
Arm 7 & Suction + irrigation & P1 to P8 & Suction pressure, pH \\
Arm 8 & Final imaging + drain placement & P1, P4, P8 (light load) & ICG fluorescence, drain ID \\
\bottomrule
\end{tabularx}
\end{table}

\begin{table}[H]
\caption{Ten representative sensor channels drawn from the 640-channel inventory
(80 channels per arm).}
\label{tab:sec06-sensors}
\centering
\begin{tabularx}{\textwidth}{L{1.0cm} L{4.6cm} L{2.2cm} L{2.4cm} Y}
\toprule
\textbf{Arm} & \textbf{Channel name} & \textbf{Rate} & \textbf{Unit} & \textbf{Dynamic range} \\
\midrule
1 & tip\_force\_x & 100 kHz & N & 0 to 3.0 N \\
1 & joint\_torque\_3 & 10 kHz & N$\cdot$m & $-50$ to $50$ \\
2 & tip\_force\_z & 100 kHz & N & 0 to 3.0 N \\
3 & retractor\_force & 10 kHz & N & 0 to 5.0 N \\
4 & icg\_fluorescence & 10 kHz & a.u. & 0 to $65{,}535$ \\
4 & nir\_depth & 10 kHz & mm & 0 to 50 \\
5 & ring\_tension & 10 kHz & N & 0 to 1.0 N \\
6 & rf\_power & 10 kHz & W & 0 to 80 \\
7 & suction\_pressure & 10 kHz & kPa & $-90$ to 0 \\
8 & drain\_id & 10 kHz & enum & 0 to 4 \\
\bottomrule
\end{tabularx}
\end{table}

\paragraph{Physical AI System Description fields (21 CFR $\S$312.23(g)).}
Because the device participates in the investigational drug delivery pathway, the
combined IND/IDE carries a full Physical AI System Description meeting the eleven
fields of 21 CFR $\S$312.23(g) \cite{kawchak2026adapt312}. The fields, and the
location of the corresponding evidence in this submission, are summarized below.

\begin{itemize}[leftmargin=1.4em,itemsep=1pt,topsep=2pt]
\item[(i)] \textit{System identification and classification:} manufacturer, model
designation, software version, surgical robot category, and the current Unified
Safety Level (USL) score computed across simulation switching, AI integration,
cross-robot sharing, and clinical-trial collaboration.
\item[(ii)] \textit{Role in drug administration:} secondary role (the device does
not itself dispense daraxonrasib; it conducts the resection and reconstruction and
generates the operative and pathology inputs to the restart advisory), with the
permitted autonomy levels and human-oversight requirements for each procedure.
\item[(iii)] \textit{Hardware specifications:} 56 degrees of freedom, 0.05 mm root
mean square positioning at $1{,}200$ mm/s, 100 kHz force sensing, the $\leq 3$ N
per-arm and $\leq 18$ N cumulative tip-force limits, and end-effector details.
\item[(iv)] \textit{Software and AI/ML specifications:} the operating and
middleware stack, the control algorithms, and the on-premises advisory LLM
together with its training data characteristics, validation methodology,
performance metrics, and Predetermined Change Control Plan (PCCP).
\item[(v)] \textit{Simulation validation data:} the Phase 0 campaign results,
including cross-framework consistency and the sim-to-real gap quantification
(trajectory gap $< 2$ mm, tip-force gap $< 0.5$ N).
\item[(vi)] \textit{Digital twin specifications:} the digital-twin methodology,
calibration sources, update frequency, and validation metrics used for operative
rehearsal and intra-operative guidance.
\item[(vii)] \textit{Safety systems:} the emergency-stop mechanisms (hardware and
software), force and torque limits, collision detection, workspace-boundary
enforcement, fail-safe behaviors, and the pre-procedure safety matrix.
\item[(viii)] \textit{Cybersecurity assessment:} the network architecture,
deny-by-default authentication and authorization, encryption, vulnerability
management, and incident response.
\item[(ix)] \textit{MCP integration:} the Model Context Protocol server topology,
conformance level, robot capability profile, and the hash-chained audit-trail
implementation under 21 CFR part 11 \cite{ecfr2026part11}.
\item[(x)] \textit{Human-robot interaction protocols:} the operator control
stations, status displays, alarm systems, handoff procedures, and the training
requirements for operating and supervising personnel.
\item[(xi)] \textit{Maintenance and calibration:} the scheduled maintenance
intervals, calibration procedures, out-of-service criteria, and re-qualification
requirements ($\S$6.2).
\end{itemize}

\subsubsection{Dosing and Administration}

\paragraph{Daraxonrasib dosing and the 3+3 escalation.}
Daraxonrasib is administered orally once daily at one of three dose levels assigned
by the 3+3 rule with sentinel staggered enrollment. The escalation begins below the
RASolute 302 reference dose to preserve a conservative first-in-human margin and
steps up to it: DL1 is 160 mg once daily, DL2 is 220 mg once daily, and DL3 is the
300 mg once-daily reference dose \cite{rasolute302}. Three participants are enrolled
at each level and observed across a 28-day dose-limiting toxicity (DLT) window. If
no DLT occurs among the first three, the cohort escalates. If one of three
experiences a DLT, the cohort expands to six; a single DLT among six permits
escalation, whereas two or more DLTs at a level fix the maximum tolerated dose
(MTD) at the preceding level. The MTD is the highest level at which fewer than two
of six DLT-evaluable participants experience a DLT, and the recommended Phase 2
dose is selected at or below the MTD. The escalation scheme is detailed in $\S$4.3.

\paragraph{Perioperative pause-and-restart advisory.}
Daraxonrasib is paused before surgery and restarted postoperatively on a day chosen
by an advisory that is LLM-bound and human-confirmed, never autonomous. The
advisory takes two inputs: the pancreaticojejunostomy ISGPS postoperative
pancreatic-fistula grade (A, B, or C) \cite{Bassi2017} and the serum daraxonrasib
trough relative to a 0.5 ng/mL threshold. It recommends a restart on
postoperative day T+7d, T+14d, or T+21d (or a continued hold). In the 32-iteration
deterministic simulation sweep, baseline serum was 0.45 ng/mL at the operative
timepoint across all iterations (below the 0.5 ng/mL trough), and the advisory
recommended T+7d in 29 of 32 iterations (Grade A fistula with sub-threshold trough),
T+14d in 3 of 32 (Grade B or delayed serum rise), and T+21d in 0 of 32. Grade C
fistula maps to T+21d or a continued hold and to the drug stop rule of $\S$7.1.
The advisory logic is shown in Figure 8.

\begin{mermaidfig}
\node[mmin]  (pre) at (4.2,0)     {Pre-operative hold\\daraxonrasib paused\\before surgery};
\node[mmstep](in1) at (0,-2.4)    {Input: PJ ISGPS\\fistula grade\\(A / B / C)};
\node[mmstep](in2) at (8.4,-2.4)  {Input: serum trough\\vs 0.5 ng/mL};
\node[mmdec] (adv) at (4.2,-2.4)  {LLM-bound\\restart advisory\\(human confirmed)};
\node[mmgoal](r7)  at (0,-5.0)    {Restart T+7d\\Grade A,\\trough $< 0.5$\\(29 of 32)};
\node[mmstep](r14) at (4.2,-5.0)  {Restart T+14d\\Grade B / delayed\\serum rise (3 of 32)};
\node[mmdark](r21) at (8.4,-5.0)  {Restart T+21d /\\hold; Grade C\\(0 of 32 in sweep)};
\draw[mmedge] (pre) -- (in1);
\draw[mmedge] (pre) -- (in2);
\draw[mmedge] (in1) -- (adv);
\draw[mmedge] (in2) -- (adv);
\draw[mmedge] (adv) -- (r7);
\draw[mmedge] (adv) -- (r14);
\draw[mmedge] (adv) -- (r21);
\end{mermaidfig}
\figcaption{Figure 8. Daraxonrasib perioperative pause-and-restart advisory. The
restart day is keyed to the pancreaticojejunostomy ISGPS fistula grade and the
serum trough relative to 0.5 ng/mL: T+7d for Grade A with a sub-threshold trough
(29 of 32 simulation iterations), T+14d for Grade B or a delayed serum rise
(3 of 32), and T+21d or continued hold for Grade C (0 of 32 in the sweep). The
advisory is bound and human-confirmed, never autonomous.}

\paragraph{Device ``administration'': operative phases and the three anastomoses.}
The device has no pharmacologic dose; its administration is the conduct of the
pancreaticoduodenectomy across eight operative phases (P1 through P8) under
continuous human oversight. Phases P1 through P4 comprise the dissection and
specimen removal (Arms 1 and 2 dissecting under Arm 3 vessel control and Arm 4
near-infrared indocyanine-green imaging); the specimen is removed at the end of
P4 and reconstruction begins. The three reconstructive anastomoses are then
constructed in sequence under closed-loop ring-tension control by the Arm 5
anastomosis probe, which issues a soft warning whenever measured ring tension
falls outside the controller band: the duct-to-mucosa pancreaticojejunostomy (PJ)
in P5 with a band of 0.30 to 0.60 N, the end-to-side hepaticojejunostomy (HJ) in
P6 with a band of 0.20 to 0.50 N, and the antecolic gastrojejunostomy (GJ) in P7
with a band of 0.40 to 0.80 N; P8 covers final imaging, hemostasis, and drain
placement. The pancreaticojejunostomy is the dominant determinant of postoperative
pancreatic-fistula risk and is therefore the anastomosis on which the fistula-grade
input to the drug advisory is based. The reconstruction sequence and the
ring-tension bands are shown in Figure 9.

\begin{mermaidfig}
\node[mmin]  (rem) at (0,0)       {Specimen removed\\(end P4)\\reconstruction begins};
\node[mmgoal](pj)  at (4.2,0)     {Pancreatico-\\jejunostomy (PJ)\\duct-to-mucosa, P5\\0.30 to 0.60 N};
\node[mmgoal](hj)  at (8.4,0)     {Hepatico-\\jejunostomy (HJ)\\end-to-side, P6\\0.20 to 0.50 N};
\node[mmgoal](gj)  at (8.4,-2.8)  {Gastro-\\jejunostomy (GJ)\\antecolic, P7\\0.40 to 0.80 N};
\node[mmstep](mon) at (4.2,-2.8)  {Ring-tension monitor\\(Arm 5)\\soft-warn outside\\band};
\draw[mmedge] (rem) -- (pj);
\draw[mmedge] (pj) -- (hj);
\draw[mmedge] (hj) -- (gj);
\draw[mmedged] (mon) -- (pj);
\draw[mmedged] (mon) -- (hj);
\draw[mmedged] (mon.east) -- ++(0.4,0) |- (gj.west);
\end{mermaidfig}
\figcaption{Figure 9. The three reconstructive anastomoses and their controller
ring-tension bands: the duct-to-mucosa pancreaticojejunostomy (0.30 to 0.60 N), the
end-to-side hepaticojejunostomy (0.20 to 0.50 N), and the antecolic
gastrojejunostomy (0.40 to 0.80 N), each guarded by the Arm 5 ring-tension monitor
that issues a soft warning outside the band.}

\subsection{Preparation, Handling, Storage, and Accountability}

\subsubsection{Acquisition and Accountability}

Daraxonrasib is supplied by the sponsor to the investigational-product pharmacy at
each participating site and is dispensed only by the site pharmacist or an
authorized designee against a written prescription tied to an enrolled, eligible
participant and a confirmed dose level. The pharmacy maintains a drug-accountability
log recording, for every unit, the lot number, the quantity received, the quantity
dispensed, the quantity returned, and the quantity destroyed, reconciled against
shipment records at each monitoring visit. Returned and unused product is retained
for reconciliation and destroyed only after sponsor authorization. The
investigational device is provided under the IDE and is not a consumable; its
accountability is established through the asset record, the per-procedure
pre-procedure safety matrix, and the hash-chained audit trail that captures every
operative use ($\S$6.4). Custody of both components is restricted to authorized,
trained personnel, and all transfers are documented.

\subsubsection{Formulation, Appearance, Packaging, and Labeling}

Daraxonrasib is provided as oral tablets in the strengths required to compose the
DL1 (160 mg), DL2 (220 mg), and DL3 (300 mg) once-daily doses. Tablets are packaged
in child-resistant, light-protective containers labeled in accordance with 21 CFR
$\S$312.6, bearing the caution statement ``Caution: New Drug. Limited by Federal
(or United States) law to investigational use,'' together with the protocol number,
lot number, storage conditions, and a unique container identifier; no information
that would unblind an outcome assessor appears on the label. Because this is an
open-label study, the participant-facing label identifies the product and the
assigned daily dose. The investigational device and its single-use accessories are
labeled per 21 CFR part 812 with the investigational-device caution statement and
are not relabeled or repackaged at the site.

\subsubsection{Product Storage and Stability}

Daraxonrasib is stored at controlled room temperature in a secured,
access-controlled, temperature-monitored location within the investigational-product
pharmacy, in its original light-protective packaging, until dispensed. Continuous
temperature monitoring with documented excursion handling is maintained; any
excursion outside the labeled range is quarantined and is not dispensed until the
sponsor confirms continued suitability against the stability data on file. Stability
is supported by the sponsor's stability program, and product is dispensed only
within its labeled expiry or retest date. The investigational device is stored and
maintained in its qualified clinical environment between procedures under the
maintenance and calibration program below.

\subsubsection{Preparation and Device Maintenance and Calibration}

\paragraph{Drug preparation.}
No compounding or reconstitution is required; preparation consists of the
pharmacist verifying the participant identity, the assigned dose level, and the
expiry, then dispensing the once-daily oral dose with administration instructions.
Tablets are swallowed whole with water and are not crushed, split, or chewed unless
a future protocol amendment provides validated instructions.

\paragraph{Device maintenance and calibration (21 CFR $\S$312.23(g)(xi)).}
Before each clinical procedure the system must pass the pre-procedure safety matrix:
verification of the 640-channel sensor stack, confirmation of the per-arm force
calibration against the $\leq 3$ N per-arm and $\leq 18$ N cumulative caps,
exercise of the 10 kHz heartbeat with its $100\,\mu$s watchdog, a successful
software and hardware E-stop self-test, and confirmation that the positioning
accuracy remains within the 0.05 mm root mean square specification at $1{,}200$
mm/s. Scheduled maintenance, periodic force and positional calibration, and
performance verification are performed at the manufacturer-specified intervals and
documented. A system that fails any pre-procedure check, exhibits an out-of-tolerance
calibration, or experiences a hard-stop E-stop event is taken out of service and is
returned to clinical use only after re-qualification against the same matrix and
sponsor authorization. All maintenance, calibration, and out-of-service events are
recorded in the hash-chained audit trail.

\subsection{Measures to Minimize Bias: Randomization and Blinding}

This is an open-label, single-arm, first-in-human study; there is no randomization
and no comparator arm, consistent with the dose-finding and early-feasibility
objectives ($\S$4.2) and with ICH E9 principles for early-phase trials. Because no
treatment assignment is masked, bias control is concentrated on the outcome
assessments most vulnerable to it. The primary surgical-efficacy endpoint, R0
resection status, is determined by a pathologist who is masked to the assigned
daraxonrasib dose level, to the intra-operative device telemetry, and to the
LLM advisory record, and who applies a standardized synoptic margin protocol;
discordant or indeterminate margins are adjudicated by an independent masked
second pathologist. Radiographic response is read centrally against RECIST by a
reviewer blinded to dose level and to clinical course, and the serious
adverse-event and Physical AI safety adjudications described in $\S$8.3 are
performed by reviewers independent of the operative team. Sentinel staggered
enrollment with case-by-case data safety monitoring board review further protects
the integrity of the safety read-out by preventing within-cohort contamination of
clinical judgment.

\subsection{Study Intervention Compliance}

\paragraph{Daraxonrasib adherence.}
Adherence is documented through the dispensing and accountability log, returned-tablet
counts at each visit, and a participant dosing diary reconciled against the
pharmacy record. Adherence is corroborated objectively by the serum daraxonrasib
assay, which both confirms exposure and supplies the trough value used by the
restart advisory ($\S$6.1.2); a documented trough discordant with the reported
intake prompts adherence counseling and is recorded in the case report form.
Material non-adherence is captured as a protocol deviation and informs the
DLT-evaluability determination ($\S$9).

\paragraph{Device compliance.}
Device compliance is enforced before and during every procedure rather than
reconstructed afterward. No operative phase may begin until the pre-procedure
safety matrix passes in full, and the 10 kHz heartbeat, the $100\,\mu$s watchdog,
the force caps, and the vascular safety-zone gate operate continuously throughout
the case. Every advisory command, every safety-gate verdict, every force and
trajectory sample at the recorded rates, and every E-stop event is written to a
21 CFR part 11 hash-chained, append-only audit trail \cite{ecfr2026part11}, so that
compliance with the cyber-physical envelope is verifiable for each case and any
deviation is detectable and time-stamped. The audit trail also satisfies the
$-24$ h to $+72$ h preservation requirement that applies around a reportable
Physical AI adverse event ($\S$8.3.6).

\subsection{Concomitant Therapy}

All medications, biologics, blood products, radiotherapy, and clinically significant
non-pharmacologic interventions taken from the time of informed consent through the
end of safety follow-up are recorded in the case report form with indication, dose,
route, and dates. Standard perioperative supportive care is permitted and
encouraged, including antiemetics, analgesia, venous thromboembolism prophylaxis,
peri-operative antimicrobial prophylaxis, proton-pump inhibitors, and
pancreatic-enzyme replacement as clinically indicated, with attention to agents that
may interact with daraxonrasib exposure. Prohibited during study participation are
other systemic anticancer therapy (cytotoxic chemotherapy, other targeted agents,
or immunotherapy) outside the protocol, other investigational agents or
investigational devices, and any concomitant use of a second surgical robotic
system for the index operation. Strong inducers or inhibitors of the principal
metabolic pathway of daraxonrasib are avoided where feasible; when an interacting
agent is clinically unavoidable, its use is documented and reviewed against the
exposure data.

\subsubsection{Rescue Medicine}

Rescue is available along two axes. For surgical and device-related complications,
the operative team may at any time convert to manual or open technique ($\S$7.1) and
deliver standard rescue care: transfusion and hemostatic resuscitation for bleeding,
vasoactive support for hemodynamic instability, percutaneous or operative drainage
and broad-spectrum antimicrobials for a clinically relevant postoperative
pancreatic fistula or intra-abdominal collection, and somatostatin-analogue therapy
where indicated for high-output fistula. For drug toxicity, daraxonrasib is held or
discontinued per the DLT and stopping rules ($\S$7.1) and managed symptomatically
with the appropriate supportive agents (for example antidiarrheals and aggressive
dermatologic care for the rash and gastrointestinal toxicities characteristic of
the class), with dose interruption and reduction governed by the protocol. All
rescue interventions are recorded as concomitant therapy and, where the criteria are
met, as adverse or serious adverse events.
